% Fragmento público generado desde propietarios íntegros.
% Residencia: 52.1.1.
% Fuente: ../incoming/radion_monografia_20260827/manuscrito/sections/02_genealogia.tex:1-180
\noindent La généalogie opératoire précède la mesure angulaire et détermine les variables conservées par chaque lecteur.\par\medskip

\subsubsection{Support APP--TRIT--TPK et construction de l'état enrichi}

\paragraph{Antériorité de l'état enrichi par rapport à la publication circulaire}

La définition scolaire du radian commence par un cercle déjà donné, un centre déjà
donné et une longueur d'arc déjà mesurable. Si le rayon et l'arc ont la
même longueur, l'angle sous-tendu mesure un radian. La définition est correcte
au sein de la géométrie euclidienne, mais ne répond pas à une question antérieure :
quelle structure permet à un tour de posséder phase, orientation, action et
mémoire, et pourquoi l'unité naturelle est-elle liée à \(2\pi\) ?

HMT inverse l'ordre causal. Elle construit d'abord l'état fini, puis sa
prolongation coinductive, ensuite le tour et seulement à la fin sa publication
angulaire. Le radion ne corrige pas le radian conventionnel ; il le relève vers son domaine
généalogique.

\begin{typebox}
\[
\APP\to\TRIT\to\TPK\to X_{\mathrm{enr}}
\to\mathcal C^{\mathrm{disc}}_{\mathrm{cont}}
\to\bigl(\pih,e_{\HMT},\varphi_{\HMT},\alphah,A,C^*,\Delta_4,\hret\bigr)
\to\mathfrak r_\sigma.
\]
Les mathématiques conventionnelles n'apparaissent qu'ensuite comme langage de réalisation :
\(S^1\), ellipse symplectique, métrique de Klein, espace de Hilbert, circuit
LC, écran de Minkowski et équations de champ.
\end{typebox}

\paragraph{Support bivarié d'APP : feuillets, résidus, quotients et retenue}

La plateforme d'arithmétique positionnelle (APP) conserve simultanément deux
réalisations des symboles \(1,\dots,9\) : un feuillet additif et un feuillet
multiplicatif. Un nombre n'est pas réduit à sa valeur publiée. Avec la valeur
subsistent feuillet, orientation, quotient, reste, retenue, frontière, chemin et
survie.

Le recensement de base comporte trois strates qui ne doivent jamais être identifiées par simple
coïncidence de cardinalité :
\[
104,976\longrightarrow468\;U_6\longrightarrow243\;w_6
\subset\F_3^6,\qquad |\F_3^6|=729.
\]
Les \(468\) émissions décimales, les \(243\) mots visibles et l'espace ambiant
de \(729\) mots sont des objets distincts. L'espace ambiant sera décisif pour la
normalisation \(20/81\).

En base cubique \(B=1000\), la soustraction orientée de deux mots de triades
s'effectue de droite à gauche :
\begin{equation}
u_i=t_i-v_i+c_{i+1},\qquad
r_i=u_i\bmod1000,\qquad
c_i=\frac{u_i-r_i}{1000}.
\label{eq:subacarreo}
\end{equation}
Les opérandes et la retenue distante étant fixés, quotient et reste sont uniques. Tel est
l'opérateur qui, plus loin, publiera les chiffres de \(\epsone\). Il n'y a
aucune décimale choisie après avoir vu le radian.

La complétion
\[
10^3=729+243+27+1
\]
n'est pas une curiosité décimale. Elle exprime la conservation de l'unité lors du passage
de l'espace ambiant ternaire à une chambre de mille positions. La différence
\(1000-729=271\) est la couronne complétante ; l'identité correcte est
\(999=729+270\), non \(729+271\).

\paragraph{Orientation ternaire et classification des 3 régimes opératoires du TRIT}

Le TRIT attribue à chaque état local l'un de trois régimes orientés. Dans la
convention temporelle active,
\[
\begin{aligned}
+1&\longleftrightarrow\text{retour, mémoire, passé},\\
0&\longleftrightarrow\text{seuil, présent},\\
-1&\longleftrightarrow\text{ouverture bidirectionnelle, futur}.
\end{aligned}
\]
La réalisation quadratique distingue :
\begin{equation}
J^2=-I,\qquad N^2=0,\qquad R^2=I.
\label{eq:tres-regimenes}
\end{equation}
Ainsi apparaissent, respectivement, la rotation elliptique, le seuil parabolique et
la séparation hyperbolique. On ne mélange pas trois métriques. On conserve trois
actions typées sur un antécédent commun.

La racine interne de \(-1\) n'est pas importée en supposant d'avance le plan
complexe. Dans le bloc APP, deux projections \(P,Q\) de rapport
\(r=3/4=90/120\) produisent
\[
J_{\mathrm{comm}}=\frac{[P,Q]}{\sqrt{r(1-r)}},\qquad J_{\mathrm{comm}}^2=-I,
\]
tandis que
\[
R_{\mathrm{comm}}=2P-I,\qquad R_{\mathrm{comm}}^2=I,\qquad
R_{\mathrm{comm}}J_{\mathrm{comm}}=-J_{\mathrm{comm}}R_{\mathrm{comm}}.
\]
Le couple génère la forme minimale de \(\mathrm{Cl}_{1,1}\simeq M_2(\R)\) :
rotation et ouverture émergent ensemble, mais restent des opérations différentes.

\paragraph{Composition dynamique du TPK : sélection, transport, mémoire et retour}

Le Topological Phase Kernel n'est pas le nom d'une boîte d'objets. C'est la
composition qui sélectionne la phase, transporte la cardinalité, relève les feuillets,
enregistre la mémoire et exécute le retour. La dépendance fondamentale est
\[
U_t\longrightarrow\Gamma_9\longrightarrow K_{\mathrm{ph}}.
\]
L'holonomie nonadique \(\Gamma_9\) agit sur l'état enrichi ; la
publication \(K_{\mathrm{ph}}\) est ultérieure et conserve une mémoire réduite. L'identité
\(K_{\mathrm{ph}}^9|_r=E_+E_\times\) ne transforme pas \(K_{\mathrm{ph}}\) en
générateur.

Le certificat nonadique conserve la chaîne complète
\[
w_6\to w_{12}\to w_{18}\to w_{24}\to w_{30}
\to R_{36}\to G_9,
\]
avec \(396\) niveaux, \(2376\) trits, \(43\) tours, \(387\) paires
\((K,K+9)\) par canal et cinq régions de \(\pi\). La phase satisfait
\[
g(k+9)=g(k),
\]
mais l'état complet ne revient pas :
\[
B_{k+9}\ne B_k.
\]
Tel est le germe mathématique de l'hélice : la projection angulaire se ferme ; la
mémoire avance.

\paragraph{Extension temporelle non scindée \(C_{12}\to C_{108}\to C_9\)}

La mémoire de douze pas s'organise au moyen de
\begin{equation}
0\longrightarrow C_{12}\longrightarrow C_{108}
\longrightarrow C_9\longrightarrow0.
\label{eq:extension-108}
\end{equation}
Le cocycle
\[
c(r,r')=\left[\left\lfloor\frac{r+r'}9\right\rfloor\right]_{12}
\]
mesure la retenue entre retour de phase et avance de mémoire. Neuf pas ferment
la phase observable \(C_9\) ; ils ne réinitialisent pas le registre dodécaphasique.

Il convient de séparer quatre objets :
\begin{center}
\begin{tabularx}{.95\textwidth}{>{\bfseries}l X}
\toprule
\(C_{12}\) & groupe cyclique de douze positions ;\\
\(\Rstate\) & registre enrichi \((E_m,\tau_m,\sigma_m,\kappa_m,\Lambda_m)_{m=1}^{12}\) ;\\
\(\Theta_{108}\) & publication brute dans le calendrier de 108 ;\\
\(\Gamma_{108}\) & histoire ou holonomie relevée.\\
\bottomrule
\end{tabularx}
\end{center}
C'est pourquoi l'expression vague \enquote{post-\(C_{12}\)} sera remplacée par
\enquote{postérieur au registre dodécaphasique \(\Rstate\), dans sa lecture
sexadique} lorsque telle sera l'application effective.

\paragraph{Émergence corrélative des 5 constructions du continu à partir de l'état unique}

La structure discrète du continu ne s'obtient pas en assemblant cinq théories
indépendantes. Ses cinq aspects internes ---spectral, solénoïdal, de jauge,
cohomologique et déterminantal--- émergent corrélativement du même état
APP--TRIT--TPK. Cet ouvrage ne démontre pas de nouveau le traité intégral, mais
préserve le fait causal dont le radion a besoin : la section coinductive qui
génère \(\pi,e,\varphi\) et \(\alpha\) ne vit pas sur une droite réelle nue ; elle vit
dans un objet doté de fibres, de relations, d'orientation, de retenue, de mémoire, de terminal et de
lecteur.

\begin{statusbox}
\textbf{Falsificateur de généalogie.} Si une formule ultérieure reçoit
\(180/\pi\), \(\epsR\), une masse ou une table métrologique pour choisir une phase,
un coefficient ou une branche, elle cesse d'être une génération HMT-first. La confrontation
externe peut reconnaître une sortie ; elle ne peut pas la sélectionner.
\end{statusbox}
